\documentclass[10pt, aps, rmp,  twocolumn, superscriptaddress, floatfix, showkeys]{revtex4-2}
\usepackage[shortlabels]{enumitem}
\bibliographystyle{plainnat}
%\usepackage[latin1]{inputenc}
\usepackage{multirow}
\usepackage{booktabs}
\usepackage{amsmath}
\usepackage{amsfonts}
\usepackage{graphicx}
\usepackage{calrsfs} 
\usepackage{times}
\usepackage[sort&compress]{natbib}
\usepackage{subfigure}
\usepackage{xcolor}
\usepackage{calrsfs}
\usepackage{soul}
\usepackage{booktabs, multirow}
\usepackage{comment}
\usepackage{verbatim}
\usepackage{framed}
\usepackage{lineno}
\linenumbers

\usepackage{xcolor}
\usepackage{mdframed}

\newmdenv[
  backgroundcolor=gray!10,
  linecolor=gray!10,
  linewidth=0pt,
  innerleftmargin=10pt,
  innerrightmargin=10pt,
  innertopmargin=10pt,
  innerbottommargin=10pt,
  skipabove=10pt,
  skipbelow=10pt
]{graybox}

% Preamble code to extract detailed section word counts
\newcommand{\detailtexcount}[1]{%
  \immediate\write18{texcount -v -merge -q #1.tex > #1.wcdet}%
  \verbatiminput{#1.wcdet}%
}

\providecommand{\keywords}[1]{\textbf{\textit{Keywords---}} #1}


\begin{document}
\newcommand{\fz}[1]{\textcolor{blue}{#1}}
\newcommand{\tbd}[1]{\textcolor{blue}{[TBD] #1}}

%\title{Alternative genome instability states in a digital model of tumorigenesis with whole genome duplication}

%\title{Evolution of critical cancer instability after whole genome duplication}
%\title{Evolutionary paths towards critical cancer instability}
%\title{Self-Organized cancer fragility emerges from chromosomal instability}
\title{A Probabilistic Ensemble-Based Bayesian Method for Mechanistic GRN Inference}
%\title{Genomically unstable evolution drives tumours to the edge of viability}

  
\providecommand{\DFA}{Department of Physics and Astronomy ``Galileo Galilei'', University of Padova, Padova, Italy}
\providecommand{\PNC}{Padova Neuroscience Center, University of Padua, Via G. Orus 2, 35131 Padova, Italy}
\providecommand{\PCNM}{Padua Center for Network Medicine, University of Padua, Via F. Marzolo 8, 315126 Padova, Italy}
\providecommand{\INFN}{Istituto Nazionale di Fisica Nucleare, Sez. Padova, Italy}

\author{Francesco Zambelli}
\affiliation{\DFA}
\author{Valeria D'Andrea}
\affiliation{\DFA}
\author{Manlio De Domenico}
\affiliation{\DFA}
\affiliation{\PNC}
\affiliation{\PCNM}
\affiliation{\INFN}

\vspace{0.4 cm}
\begin{abstract}
\vspace{0.2 cm}  
Recovering the interaction structure of a complex stochastic dynamical system from time-series observations is an inherently ill-posed problem: the posterior over network configurations is high-dimensional, multimodal, and shaped by the non-equilibrium nature of the available data. We propose a Bayesian ensemble inference framework grounded in the Onsager-Machlup path probability formalism, which provides a rigorous likelihood for SDE-governed trajectories without requiring stationarity, and is coupled with a node-wise block-coordinate Metropolis-Hastings sampler that exploits the factorization of the path probability to make inference tractable. A evidence-estimation-based model comparison scheme jointly recovers the dynamical parameters alongside the network topology. Applying the framework to gene regulatory network (GRN) inference, we demonstrate on synthetic data that ensemble-based posterior analysis produce more reliable results, in particular when considering perturbation predictions and network topological statistics, than single-point estimation. The framework is general and can be applied to any mechanistic SDE model, offering a powerful tool for uncertainty-aware inference in complex stochastic systems.
\vspace{0.2 cm}
\end{abstract}

\keywords{Cancer, Tumor Growth, Genomic Instability, Error Threshold}

\maketitle

% network inference is important
Inferring the structure of a complex dynamical system from time-series observations is a fundamental challenge in the physical, biological, and data sciences~\cite{newman2018network,peixoto2025network,nitzan2017revealing,timme2007revealing, brunton2016discovering}. 
%%% \valeria{aggiungerei altre referenze, ad esempio (non necessariamente in quest'ordine):}
%%% Aggiunti: Peixoto PRX 2025, Newman NatPhys 2018, Barzel NatPhys 2013, Peel NatComm 2022, Nitzan SciAdv 2017, Timme 2007
%%% Guimerà PNAS 2009: data are referred to edges

% in particular GRN inference is important
In particular, inferring the structure of gene regulatory networks (GRNs) from gene expression is a crucial problem in systems biology, with important implications for understanding cellular function and disease mechanisms~\cite{karlebach2008modelling, barabasi2011network}. Knowing the structure of the GRN allows us to predict how the system will respond to perturbations such as gene knockouts or drug treatments, and to identify potential therapeutic targets. 
% structure and dynamics are intertwined
However, structure and dynamics are deeply intertwined in complex systems~\cite{gross2008adaptive,barzel2013universality,ghavasieh2024diversity,strogatz2024nonlinear,barzon2024unraveling,sole2001complexity,de2026decoding}, and thus inferring the network structure without explicitly modeling the underlying dynamics can lead to inaccurate and unreliable results. %chiedere a manlio per la referenza
% paramter non-identifiability and structural non-identifiability
When the system is governed by nonlinear, stochastic dynamics, and the observations are noisy and limited in time, the problem becomes particularly severe: multiple distinct configurations of the interactions can produce trajectories that are statistically indistinguishable from the data~\cite{villaverde2014reverse}. These difficulties, broadly referred to as structural non-identifiability, call for inference approaches that go beyond point estimation, instead quantifying and propagating uncertainty over the full space of plausible model configurations.
% likelihood discontinuity and network emergence
Another important factor is that when dealing with networks, even small changes in the parameters can lead to large changes in the system's behavior, due to the emergence of collective phenomena and critical transitions~\cite{boccaletti2006complex,dorogovtsev2008critical,pastor2015epidemic, strogatz2024nonlinear}. This leads to highly discontinuous likelihood landscapes, where the maxima are isolated and separated by low-probability barriers. For this reason, in most of the cases, standard incremental optimization methods are not good solutions for such problems.
%%% \valeria{altri paper eventalmente da citare: , Pastor Satorras RMP 2015,...}
%%%Aggiunti: Boccaletti Phys Rep 2006,  Dorogovtsev 2008 RMP
% importance of bayesian inference and uncertainty quantification
A natural and principled route to this goal is Bayesian inference, in which a posterior distribution over model parameters is obtained by combining a likelihood function for the observed data with prior knowledge about the system~\cite{peixoto2019network, newman2018network,raimondo2021measuring,peel2022statistical,peixoto2025network}. This approach allows us to quantify uncertainty over the inferred interactions, which is crucial for guiding experimental validation and for making reliable predictions about the system's behavior.

% existing methods
%% correlation based methods
Most of the methodologies proposed to tackle this problem rely on statistical associations between observed variables, such as correlation or mutual information, to infer the presence of interactions between them~\cite{margolin2006aracne,marbach2012wisdom, marku2023time}. Additionally, they do not provide measures of uncertainty over the inferred interactions, which is crucial for guiding experimental validation and for making reliable predictions about the system's behavior. 
%% bayesian structure learning and graphical models, DAGs
Other methods attempt to learn the posterior distribution over network configurations using Bayesian structure learning techniques, often based on graphical models and directed acyclic graphs (DAGs)~\cite{glymour2019review}. These methods can provide uncertainty estimates, but they typically assume that the underlying system is governed by a DAG structure, which is not always valid in real systems, and thus they cannot properly capture cyclic interactions, which are common in biological systems. Additionally, they often rely on assumptions of stationarity and equilibrium, which are not always valid in real systems, and thus they cannot fully exploit the information contained in non-equilibrium trajectories, which are often the only ones available in real experiments.
%% cyclic structure learning methods
Given the inherently cyclic nature of gene regulatory systems, it is crucial to consider methods that can capture such interactions. Multiple approaches have been proposed to solve the issue in different fields~\cite{huang2020causal,runge2019detecting}. However, these methods have limitations, such as single point estimation, lack of uncertainty quantification, and reliance on assumptions that may not hold in real systems, such as stationarity and equilibrium.
%% work atanokvic
A recent approach~\cite{atanackovic2023dyngfn} addresses some of these limitations by proposing a Bayesian framework for inferring cyclic interactions from time-series data by modeling the system's latent drift function through either linear solvers or non-linear MLPs. While this method represents a significant step forward, it still does not explicitly define the mechanistic SDE model underlying the dynamics, and thus remains of difficult interpretation, comparison with other mechanistic models of gene regulation and with the risk of data overfitting. Additionally, it does not provide a principled way to select the dynamical parameters of the SDE model, such as noise amplitude and relaxation rates.

% Our proposal  
%% Building likelihood
What is needed to overcome these limitations is a mechanistic, ensemble-based Bayesian framework that can rigorously define a likelihood for non-equilibrium trajectories generated by a mechanistic SDE model, while also providing a tractable sampling algorithm to explore the high-dimensional and multimodal posterior distribution over network configurations.
%% likelihood based on path probability
We propose a rigorous method to build a likelihood for non-equilibrium trajectories generated by a mechanistic SDE model, based on the Onsager-Machlup path probability formalism~\cite{onsager1953fluctuations}. This path-probability likelihood has a decisive advantage over approaches based on stationary distributions: it is valid for both transient and non-equilibrium trajectories, enabling the full information content of time-series data to be exploited without requiring prior knowledge of the system's stationary state or equilibrium properties.
%% explicit prior integration to avoid non-identifiability and overfitting
Data alone is often insufficient to identify the true network structure, due to the fundamental non-identifiability of the problem. To overcome this issue, we explicitly integrate prior knowledge about the system into the inference process, by defining a prior distribution over network configurations that encodes relevant constraints based on biological knowledge. This allows us to guide the inference process towards more plausible configurations, and to avoid overfitting, which is a common problem in high-dimensional inference problems.
%% multimodality and parallel tempering
High probability regions in the parameter space are often separated by low-probability barriers, leading to a multimodal posterior distribution. To efficiently explore this complex landscape, we use a node-wise block-coordinate Metropolis-Hastings sampling algorithm that exploits the factorization of the path probability over target nodes to make posterior sampling tractable, while maintaining ergodicity guarantees. We combine it with a parallel tempering scheme to further enhance the exploration of the parameter space and to being able to escape from local minima.
%% model comparison
%% final output
What we obtain from this procedure is an ensemble of configurations for the parameters, each weighted by its posterior probability, that can be analyzed to extract probabilistic measures on the system under analysis rather than a single point estimate, which is crucial for guiding experimental validation and for making reliable predictions about the system's behavior~\cite{peel2022statistical}.

% Applications
%% Syntetic network on nonlinear SDE model
As a primary application, we start from a nonlinear SDE model with sigmoidal activation functions -- capturing saturation, oscillatory, and noise-driven regimes -- and then we perform inference on synthetic networks. We show that the ensemble approach provides a more reliable reconstruction of the true network topology than single-point estimation, in particular for the prediction of gene knockout effects. We also demonstrate how incorporating sparse prior knowledge about network connectivity and known regulatory interactions improves inference in data-limited regimes.
%% RNA velocity data
%% Generalizability
%%%\valeria{l'ultima frase è più per Conclusions}

%%%\valeria{Attenzione che i pannelli non sono labeled con le lettere come in caption. Alcune scritte sono troppo piccole (es. la legend della sparsity o le scritte dentro i singoli nodi, meglio cancellarle o ridimensionarle. Infine nel pannello centrale: immagino che le piccole matrici in basso a dx si riferiscano a misure e errore per diversi tipi di traiettorie: inserire la spiegazione in caption)}
\begin{figure*}
    \centering
    \includegraphics[width=\linewidth]{images/Fig_1.pdf}
    \caption{\textbf{Inference procedure workflow.} In the left panel, we show the network topology and the model used to generate the synthetic data. In the central panel, we show an example of trajectory generated from the model, and the corresponding dataset of pairs $(\mathbf{x}_k, \mathbf{dx}_k)$ of states and infinitesimal increments $\mathbf{dx}_k = \frac{\Delta \mathbf{x}_k}{\Delta t}$, where $\Delta \mathbf{x}_k = \mathbf{x}_k - \mathbf{x}_{k-1}$, that can be used to perform inference on the parameters of the model. In particular we distinguish between transient and stationary phases in the temporal evolution of the system, and we show the possible different behaviors that can be observed in the trajectories, such as oscillations, saturation and low signal. In the right panel, we show a schematic representation of the inference procedure. First, we perform a node-wise analysis of the likelihood landscape to find an ensemble of configurations for the incoming edges of each node that are compatible with the observed data. Then, we sample from the configurations of the incoming edges of each node to build an ensemble of configurations for the whole network by comparing the emerging global network properties with the prior knowledge about the system. The output of the procedure is an ensemble of configurations for the parameters, each weighted by its posterior probability, that can be analyzed to extract probabilistic measures on the system under analysis.}
\label{fig:inference_workflow}
\end{figure*}

\section{Inference Problem Setting}

% general setting definition
Consider a system composed of $N$ interacting entities, such as genes in a gene regulatory network, whose state at time $t$ is described by a vector $\mathbf{x}(t) \in \mathbb{R}^N$, where $x_i(t)$ that could represent the concentration level of the transcript corresponding to gene $i$ at time $t$. 
We assume the dynamics of the system is governed by a set of stochastic differential equations (SDEs) that describe how the state of each entity evolves over time $\frac{dx_i(t)}{dt} = f(x_i(t)) + g(x_i(t))\xi(t)$, composed by a deterministic drift term that captures the interactions between the entities $f$, and a stochastic noise term that captures the intrinsic and extrinsic fluctuations in the system, with $\xi(t)$ being a standard Wiener process. We restrict the system to be subject to uncorrelated noise, and thus we can write $g(x_i(t))=D$, where $D$ is the noise amplitude.
% choise of sigmoid model for gene regulation, approximation of interaction strenghts
There are many possible choices for the specific form of the SDEs~\cite{karlebach2008modelling,vijesh2013modeling}, depending on the level of detail we want to capture and the assumptions we make about the underlying biological processes. One of the simplest and most effective models is based on the assumption that the production of each transcript is regulated by a saturating function of the concentrations of the other transcripts, while the degradation is modeled as a linear term~\cite{yu2025discover,aguirre2025gene}. Such a model captures the essential features of gene regulation, while being computationally tractable for inference. The mathematical form of the SDEs can be written as:

\begin{equation}
    \frac{dx_i(t)}{dt}=-\gamma x_i(t)+V\cdot\;\sigma\Big( \sum_j J \, A_{i,j}\,x_j\Big) + D \xi(t)
    \label{eq:ode_sigm}
\end{equation}

Where $\gamma$ is the degradation rate of transcript $i$, $V$ is the maximum production where $\sigma(x)=\frac{1}{1+e^{-x}}$ is a sigmoid function that captures the saturating nature of gene regulation, and $D \xi(t)$ represents the stochastic noise term. The regulatory interactions are defined by $J\,A_{i,j}$, with $J$ being a scaling factor and $A_{i,j}$ indicating the type of interaction between gene $j$ and gene $i$: in particular $A_{i,j}\neq 0$ indicates that node $j$ regulates node $i$, and the sign of $A_{i,j}$ indicates whether the regulation is activating or inhibitory. For simplicity, we restrict the possible values of $A_{i,j}$ to three states: $-1$ for inhibitory interactions, $0$ for absent interactions, and $+1$ for activating interactions. This is a first-order approximation that allows us to test the inference procedure in a more controlled setting, while still capturing the essential features of the dynamics and some key features of real GRNs. 

We suppose that we have temporally ordered measurements dataset $\mathbb{D}$ sampled at discrete time intervals $\Delta t$, representing the data acquisition process in real experiments $\{\mathbf{x}_k\}_{k=1}^M$. From this, we can build a dataset of pairs $(\mathbf{x}_k, \mathbf{dx}_k)$ of states and infinitesimal increments $\mathbf{dx}_k = \frac{\Delta \mathbf{x}_k}{\Delta t}$, where $\Delta \mathbf{x}_k = \mathbf{x}_k - \mathbf{x}_{k-1}$, that can be used to perform inference on the parameters of the model.
From the experimental data, we want to infer the parameters of the model, in particular the presence and the sign of the interactions between the nodes, which are encoded in the $A_{i,j}$ matrix. A robust and principled way to do this is to build a likelihood function for the observed data given the model parameters. This can be done by using the Onsager-Machlup formalism, which allows us to express the path probability of a stochastic process governed by an SDE:

\begin{widetext}
    \begin{equation}
        P(\{\mathbf{x}_k\}_{k=1}^M|\Theta)=\Big(\frac{1}{\sqrt{2\pi D\Delta t}}\Big)^M\exp\Big[-\frac{\Delta t}{2D}\sum_{k=1}^M\Big(\mathbf{dx}_k-\mathbf{f}(\mathbf{x}_{k-1}|\Theta)\Big)^2\Big]+O(\Delta t^2)
        \label{eq:pathprob}
    \end{equation}
\end{widetext}

% consideration on the likelihood definition
Some important points to consider about this likelihood definition are the followings. First, it is valid for non-equilibrium trajectories, and thus it allows us to fully exploit the information contained in the time-series data, without requiring prior knowledge of the system's stationary state or equilibrium properties. Second, it is based on a discretization of the continuous-time dynamics, and thus it has to be considered in the limit of small time steps $\Delta t$. Further details about the limits of validity of this likelihood definition are reported in the Materials and Methods section. Third, it assumes that the noise is uncorrelated and Gaussian, which is a common assumption in many biological systems, but it can be relaxed if needed by considering more general noise models. Such expression can be interpreted as constraining the increments of the state variables to be distributed around the drift term predicted by the model following the noise distribution, which in this case is Gaussian. Mathematically, this can be expressed as:

\begin{equation}
    \Delta \mathbf{x}_k \sim \mathcal{N}\Big(\mathbf{f}(\mathbf{x}_{k-1}|\Theta)\Delta t, D\Delta t\Big)
\end{equation}

% prior contribution to avoid non-identifiability and overfitting
Data alone is often insufficient to identify the true network structure, due to the fundamental non-identifiability of the problem. To overcome this issue, we explicitly integrate prior knowledge about the system into the inference process, by defining a prior distribution $\pi(\Theta)$ over network configurations that encodes relevant constraints based on biological knowledge. This allows us to guide the inference process towards more plausible configurations, and to avoid overfitting, which is a common problem in high-dimensional inference problems.

By combining the prior information with the likelihood of observing the data given the model parameters, we can obtain a posterior distribution over the parameters using Bayes' theorem:

\begin{equation}
    P(\Theta|\{\mathbf{x}_k\}_{k=1}^M) = \frac{P(\{\mathbf{x}_k\}_{k=1}^M|\Theta) \pi(\Theta)}{\int P(\{\mathbf{x}_k\}_{k=1}^M|\Theta') \pi(\Theta') d\Theta'}
\end{equation}

Once the posterior distribution is defined, we can proceed to explore the parameter space to identify the configurations of interactions that best describe the observed data, while also being consistent with our prior knowledge about the system. 

We divide the inference problem in two steps: in the first one we exploit the factorization of the path probability to perform a node-wise analysis, finding an ensemble of configurations for the incoming weights of each node that are compatible with the observed data. In the second step, we sample from the configurations of the incoming edges of each node to build an ensemble of configurations for the whole network by comparing the emerging global network properties with the prior knowledge about the system. 


\section{Mapping the Manifold of Data-Equivalent Network Weights}

The first phase of the inference is based on the the fact that the path probability factorizes over target nodes when the noise is uncorrelated, as shown in Eq.~\ref{eq:pathprob}. It is easy to verify:

\begin{equation}
    \log P(\{\mathbf{x}_k\}_{k=1}^M|\Theta')=\sum_{i=1}^N \log P_i(\{\mathbf{x}_k\}_{k=1}^M|\Theta_i)
    \label{eq:logloss}
\end{equation}

Where $\Theta_i$ is the set of parameters associated to node $i$, i.e. the weights of the edges incoming to it, that link the gene to the ones regulating it, and $P_i$ is the path probability associated to node $i$.

In this way, we can separately analyze the incoming edges of each node, and find an ensemble of configurations for them that are compatible with the observed data. In this phase we do not consider the prior contribution, and we focus on finding the configurations that maximize the likelihood, which correspond to the ones that best describe the observed data. By doing this for each node, we can build an ensemble of configurations for the whole network, that can be analyzed to extract probabilistic measures on the system under analysis.
% sampling scheme for proposing new configurations
As mentioned before, sampling from the space of network weights configurations is a challenging problem, given that small changes in the parameters can lead to large changes in the system's behavior, due to the propagation of the effects of the interactions through the network. This leads to highly discontinuous likelihood landscapes, where in most of the cases, is not possible to reach the optimal configuration by determining "one correct edge at a time", but it is necessary to find the right combination of edges to reach the optimal configuration. For this reason, in most of the cases, standard incremental optimization methods are not good solutions for such problems, in particular gradient based methods, that are not able to take into account effect of combined perturbations of the parameters. Given the discrete nature of the problem, if the number of nodes is not too large, it is possible to exhaustively explore the space of configurations for the incoming edges of each node, and to find the ones that are compatible with the observed data.
% if not possible to explore all the configurations, use parallel tempering
If such an approach is not feasible, we can use a parallel tempering scheme~\cite{hukushima1996exchange}, where the hot chains explore the parameter space more freely, while the cold chains focus on the high-probability regions, and by allowing them to exchange configurations, we can effectively explore the parameter space and find the optimal configurations for each node.
% configuration ensemble determination
Once we have a list of configurations for the incoming edges of each node, each weighted for the value of the likelihood, we can choose the ones to keep as the most representative of the likelihood landscape. In particular, such a choice is conditioned on the choice of the noise level $D$, that is a global multiplicative factor in the likelihood. In particular, by increasing the value of $D$, we can include more configurations in the ensemble, while by decreasing it, we can focus on a smaller number of configurations that are closer to the optimal one.

\section{Sampling on Emergent Network Properties}
 
% priors
% some properties emerge only by observing the network as a whole
To further refine the inference, we can consider that some properties of the system emerge only by observing the network as a whole, and thus they cannot be captured by analyzing the incoming edges of each node separately. By combining the configurations for the incoming edges of each node, we can build an ensemble of configurations for the whole network: from each of them, we can compute some global quantities, and by comparing them with the prior knowledge about the system, we can further refine the inference and find the configurations that are more consistent with our prior knowledge.
%priors definition
We chose to consider two sources of prior information: the sparsity of the network, and the presence of feedback loops. The first one is based on the fact that in general, real GRNs are sparse, and thus we can use this information to guide the inference process towards more plausible configurations, and to avoid overfitting. The second one is based on the fact that feedback loops are fundamental motifs in gene regulatory networks, and thus we expect them to be enriched in the system under analysis. 
%% contribution and temperature paramters
The contribution of the prior information can be modulated by a temperature parameter $\beta$, that allows us to control the relative weight of the likelihood and the prior contributions in the posterior distribution. By adjusting the value of $\beta$, we can effectively control the influence of the prior knowledge on the inference process, and thus we can find a balance between fitting the data and respecting our prior knowledge about the system. The posterior distribution can be written as:

\begin{equation}
    P(\Theta|\{\mathbf{x}_k\}_{k=1}^M) \propto P(\{\mathbf{x}_k\}_{k=1}^M|\Theta) \pi(\Theta)^\beta
\end{equation}

%MH-MCMC sampling scheme
By building the posterior distribution in this way, we can proceed by sampling it using a Metropolis-Hastings MCMC algorithm. To make a smoother likelihood landscape, we can consider the probability of transition between two configurations $\Theta_i$ and $\Theta_i'$ as:

\begin{equation}
    Q(\Theta_i \rightarrow \Theta_i') = \min\left(1, \frac{P(\{\mathbf{x}_k\}|\Theta_i')}{P(\{\mathbf{x}_k\}|\Theta_i)} \right)
\end{equation}

The acceptance probability for the proposed configuration can be written as:

\begin{equation}
    A(\Theta \rightarrow \Theta') = \min\left(1, \frac{P(\{\mathbf{x}_k\}|\Theta') P(\Theta') }{P(\{\mathbf{x}_k\}|\Theta) P(\Theta)}\right)
\end{equation}

with guaranteed ergodicity.

The output of the MCMC sampling procedure is an ensemble of configurations for the parameters, each weighted by its posterior probability, that can be analyzed to extract probabilistic measures on the system under analysis.

%%%\valeria{Da rivedere la caption}
\begin{figure*}
    \centering
    \includegraphics[width=\textwidth]{images/Fig_2.pdf}
    \caption{\textbf{Results of the parameter exploration.} In (a) we show the probabilities for each edge to be activating or inhibitory, as obtained from the inferred network ensemble, in the cases of no prior (top row) and sparsity prior only (bottom row). In the panel (b) we show on the left side, how the time interval between the measurements affects the inference results. It is possible to see that, as the time interval increases, the inference results become less accurate, given that the path probability is defined in the limit of small time steps, and thus it is not able to capture the dynamics of the system when the time interval between the measurements is too large. On the right side, we show the comparison of the inference results obtained by considering only the likelihood contribution or the whole posterior distribution, by plotting the ROC curve along with the area under the curve (AUC) for each of the two cases. It is clear that adding the prior contribution to the inference process allows us to improve the inference results, in particular by reducing the number of false positives. In (c) we show the parameter grid search results for the $\gamma$ and $J$ parameters, which shows a region of the parameter space in which the model evidence reaches a maximum, indicating the best combination of parameters for describing the data. The true values of the parameters are indicated by the white cross.}
    \label{fig:fig2}
\end{figure*}

\subsection{Model-Evidence Based Coupling Parameter Selection}

To make the method more generalizable to different systems, we can consider the parameters of the dynamical model, such as degradation rate $\gamma$, production rate $V$, and interaction strength scaling factor $J$, as free parameters that can be selected based on the data. This is important because in real systems, such parameters are often unknown or difficult to measure, and thus it is crucial to have a method that can infer them from the data, alongside with the network structure.
%joint learning is unfeasible
Jointly learning the network structure and the dynamical parameters is a challenging problem at this stage, because even small changes in the parameters can lead to completely different edge configurations to be compatible with the observed data, and thus it is not possible to use a standard optimization procedure to find the optimal combination of parameters.
% best posterior configuration
What we want to get is a model comparison framework that allows us to compare different choices of the parameters, and to select the one that best describes the data. To do this, we can perform a grid search over the parameter space, and for each combination of parameters, we find the best posterior configuration of the network weights. This is allowed by the fact that, with the trajectory transformation procedure described in the Materials and Methods section, the parameters range can be effectively reduced, and thus the grid search can be performed in a reasonable time.

This allows us to apply the model to a potentially wide range of systems, without the need to set the parameters a priori, and thus making the method more generalizable and applicable to different biological contexts.

\begin{figure*}
    \centering
    \includegraphics[width=\textwidth]{images/Fig_3.pdf}
    \caption{\textbf{Prediction of stationary state distribution of node activity after perturbation.} In (a) we show the distribution of the each node activity at the stationary state after a perturbation, as obtained from the true network (red), the best inferred network without prior (yellow), the best inferred network with sparsity prior (green), and the full ensemble of inferred networks with sparsity prior (blue). The prediction from the ensemble of inferred networks is in general more accurate, and even when the main signal does not correspond to the true outcome, there is still a lower intensity signal capturing the true value, which is not the case for the single best-scoring network. In (b) we show how the true and the inferred networks change after the perturbation, and for the network derived from the ensemble, the width of the edges is proportional to the probability of the edge to be present in the ensemble. In (c) we show the velues of $1- JS$ with $JS$ the Jensen-Shannon distance between the true and the inferred activity distributions for each node. It shows that the ensemble prediction is in general more accurate than the single best-scoring network.}
    \label{fig:perturb}
\end{figure*}

%%% \valeria{non è chiarissimo qui: stai plottando (1-JS)? Perchè va specificato anche in caption e testo, altrimenti non si capisce perchè la distribuzione con Ensemble è la migliore}

\section{Results}

% definition of synthetic data generation procedure
We test our inference framework on synthetic data generated from the dynamical model described in Eq.~\ref{eq:ode_sigm}, with a specific choice of the parameters and the network topology. For sake of simplicity, and to show the potential of the method in a more clear way, we consider a network of $N=10$ nodes, and we choose the parameters in such a way to generate oscillatory, saturated and low signal trajectories, encompassing a wide range of possible behaviors. As reported before, the weights of the edges are restricted to three states: inhibitory, absent, or activating, and the noise level and degradation rate are the same for all the nodes and are considered to be known parameters during the inference. The specific choice of the parameters and the network topology is reported in the Materials and Methods section. This is a first-order approximation that allows us to test the inference procedure in a more controlled setting, while still capturing the essential features of the dynamics and some key features of real GRNs.
% measuring
Starting from this choice of parameters and topology, we simulate the system for a time period long enough to reach the stationary state. From this, we extract a time window long enough to contain a complete oscillation of the system, similarly to what can be observed in real experiments. This can be done either in the transient or in the stationary state, given that the path probability allows us to take into account both of them, and thus to extract as much information as possible from the data.
From such trajectory we extract a set of measurements at discrete time points, simulating the data acquisition process in real experiments, and we add experimental noise to the observed concentrations.
% normalization
Then, we normalize the trajectories as described in the Materials and Methods section, and we build the dataset of pairs $(\mathbf{x}_k, \mathbf{dx}_k)$ of states and infinitesimal increments $\mathbf{dx}_k = \frac{\Delta \mathbf{x}_k}{\Delta t}$, where $\Delta \mathbf{x}_k = \mathbf{x}_k - \mathbf{x}_{k-1}$, that can be used to perform inference on the parameters of the model. 
An example of synthetically generated data is shown in Fig.~\ref{fig:inference_workflow}.
% results
We apply the inference procedure described in the previous sections to the synthetically generated data, and we analyze the results to evaluate the performance of the method.
%% heatmaps, first and second inference phase
In Fig.~\ref{fig:fig2} a.1 we show the probabilities for each edge to be activating or inhibitory, as obtained from node-wise analysis of the likelihood landscape, without considering the prior contribution. In Fig.~\ref{fig:fig2} a.2 we show the same probabilities, but obtained from the full posterior distribution, that takes into account also the prior contribution. It is evident that in the second case, the probability is more concentrated around a smaller number of edges, given that the prior contribution favors low-connectivity configurations, and thus it helps to avoid overfitting problems. Quantitatively, this can be measured by the area under the ROC curve, reported in Fig.~\ref{fig:fig2}.b (right panel), that is around to 0.8 for the first case, and increases up to around 0.88 for the second one, showing a significant improvement of the inference accuracy by including the prior contribution.
% increment of delta t
In Fig.~\ref{fig:fig2}.b (left) we show how the change in the interval between measures $\Delta t$ affects the inference results. We can see that the inference accuracy measured by the AUC, drops significantly when $\Delta t$ is increased, given that the likelihood definition is based on the assumption of small time steps, and thus it becomes less accurate when this assumption is violated. However, it is possible to see that there is a range of values for $\Delta t$ in which the inference accuracy is still relatively high, indicating that the method can tolerate some degree of violation of the small time step assumption.
% grid search for parameter selection
%%%\valeria{Controllare descrizione pannelli, che non matcha con quello mostrato in Fig2.}
In Fig.~\ref{fig:fig2}.c, we show the parameter grid search results for the $\gamma$ and $J$ parameters. It is possible to observe a region of the parameter space in which the model evidence reaches a maximum, indicating the best combination of parameters for describing the data, which is close to the true values used to generate the data.


\subsection{Perturbations as Probabilistic Predictions}

One of the main goals of gene regulatory network inference is to 
predict the effect of perturbations, such as gene knockouts. 
% how we simulate gene knockout
We simulate such a scenario by removing all outgoing edges from a specific gene in the true network, and simulating the dynamics of the perturbed system using the same SDEs as before, but with the modified network structure. We then compute the node activity distributions at the stationary state, which can be used to quantify the effect of the perturbation on the system. We do this on the true network, on the network obtained by combining the highest likelihood configuration for each node, the best-scoring network obtained from the posterior distribution, and the ensemble of inferred networks obtained from the posterior distribution (Fig.~\ref{fig:fig2}.a). 
%results
The results are shown in Fig.~\ref{fig:perturb}.b, where the node activity distributions at the stationary state in the different cases are shown, and in Fig.~\ref{fig:perturb}.c, where the Jensen-Shannon distance between the true and the inferred node 
activity distributions is shown. 
% comparison of node activity distributions, good accuracy and wider range of predictions from the ensemble
What emerges is that the prediction from the ensemble of inferred networks is in general more accurate than the single best-scoring network, but more importantly, even when the main signal does not correspond to the true outcome, there is still a lower intensity signal capturing the true value, which is not the case for the single best-scoring network. This highlights the importance of considering the full ensemble of inferred configurations, which can capture a wide range of system behaviors, that do not always correspond to the stronger signal.


\section{Discussion}
% summary of the method
In this work we propose a novel method to tackle the problem of network structure inference from data. It is based on coupling the topology of the systems interaction with a specific dynamical model, chosen to capture the essential features of the system under consideration, and to define a Bayesian inference framework to explore the parameter space and identify the configurations that best describe the observed data, while also being consistent with the prior knowledge that could help the identification of the set of parameters that best describe the system. The key element is to define a mechanistic model that predicts the true real increments between time points with a probability that is consistent with the noise of the system. This allows to avoid overfitting and to properly treat stochastic systems, allowing also to tackle the problem of non-identifiability of the parameters, which is a common problem in this kind of inference problems, by considering the whole ensemble of inferred configurations rather than relying on a single point estimate.
% use case
We use as a first application case the problem of gene regulatory network inference, which is a fundamental problem in systems biology, and we show how the proposed method can be applied to synthetically generated data, and to real data, to infer the underlying network structure and to predict the effect of perturbations. Such a method can be applied to a wide range of systems, by choosing the appropriate dynamical model and the prior knowledge that is relevant for the system under consideration, making it a powerful tool for understanding complex systems in general. 
% limitations
This work must be considered as a proof of concept, and thus it has some limitations that can be addressed in future work. 
% model choice
First, it was chosen a simple dynamical model to generate and describe the data. Processes of gene regulation can be more complex than the one captured by the model used in this work, and thus it would be interesting to extend the method to more complex models that can capture additional features of gene regulation, such time delays, distinction between mRNA and protein levels, or more complex regulatory functions. 
% noise model
Second, the method relies on the assumption of uncorrelated noise, which may not always be valid in real systems, and thus it would be interesting to extend the method to more general noise models. 
% interaction strenghts discretization
Third, the method is based on a discretization of the interaction strengths, which is a first-order approximation that allows us to test the inference procedure in a more controlled setting. However, it would be interesting to extend the method to continuous interaction strengths, which would allow us to capture a wider range of possible interactions and to make the method more generalizable. 
% model comparison
Finally, to meaningfully compare different SDE models on a given dataset which generative process is unknown, it would be necessary to extend the model comparison framework, to combine the descriptive ability of each model with the complexity of the model itself, in order to avoid overfitting and to select the model that best captures the underlying generative process of the data. This can be done by leveraging model evidence estimation techniques, such as thermodynamic integration or nested sampling~\cite{skilling2006nested, buchner2023nested}, that allow us to estimate the model evidence for each model, and thus to compare them in a principled way.
%%% \valeria{add refs}
% conclusions
In conclusion, this work presents a novel method for network structure inference from data, that combines mechanistic modeling with Bayesian inference, and that can be applied to a wide range of systems. The method is based on defining a dynamical model that captures the essential features of the system under consideration, and on using the path probability of the observed data given the model parameters to perform inference on the network structure. By considering the full ensemble of inferred configurations, rather than relying on a single point estimate, we can capture a wide range of system behaviors, and thus make more accurate predictions about the system's response to perturbations.

The proposed framework is not specific to gene regulation: any system for which an SDE model can be formulated—from neural circuits to ecological networks and epidemic models—can be treated within the same apparatus, making this a general-purpose tool for mechanistic inference in complex stochastic systems.



\section{Materials and Methods}

\subsection{Validity of the temporal discretisation.}

The Onsager–Machlup likelihood in Eq.~\eqref{eq:pathprob} relies on approximating the drift $F_i(\mathbf{x})$ by its value at the left endpoint of each interval $[t_k,t_{k+1}]$ (Euler–Maruyama freezing).
The resulting temporal error is
\[
  \varepsilon_i
  = \int_{t_k}^{t_k+\Delta t}\!\bigl[F_i(\mathbf{x}(s))-F_i(\mathbf{x}(t_k))\bigr]\,ds.
\]

Because $F_i$ is Lipschitz-continuous with constant $L_i = \gamma_i + \sum_j |w_{ij}|L_\phi$, and state fluctuations scale as $\mathbb{E}[|\mathbf{x}(s)-\mathbf{x}(t_k)|]=O(\sqrt{s-t_k})$, the unconditional bound is $\mathbb{E}|\varepsilon_i^{\text{time}}|\le\tfrac{2}{3}C_i(\Delta t)^{3/2}$.
Conditioning on the observed state $\mathbf{x}(t_k)$, a Taylor expansion of the drift gives
\[
  \mathbb{E}\bigl[\varepsilon_i^{\text{time}}\mid\mathbf{x}(t_k)\bigr] = O(\Delta t^2),
  \qquad
  \mathrm{Var}\bigl(\varepsilon_i^{\text{time}}\bigr) = O(\Delta t^3),
\]

so the conditional bias is quadratic in $\Delta t$ while the process noise is $O(\sqrt{\Delta t})$, confirming that the approximation is negligible in the likelihood for small time steps.

The practical constraint on $\Delta t$ becomes explicit when the system exhibits oscillatory dynamics.
Near a limit cycle the dominant eigenvalues of the Jacobian $J_{ij}=\partial F_i/\partial x_j$ are complex, $\lambda_\pm=\alpha\pm i\omega$, with oscillation period $T=2\pi/\omega$.
Locally, $x_j(s)-x_j(t_k)\sim A\omega(s-t_k)$, so the drift deviation scales as $F_i(\mathbf{x}(s))-F_i(\mathbf{x}(t_k))\sim L_i A\omega(s-t_k)$, yielding
\[
  \varepsilon_i
  \sim \frac{1}{2}L_i A\,\omega\,(\Delta t)^2
  = \frac{\pi L_i A}{T}(\Delta t)^2.
\] 

The dimensionless control parameter governing approximation validity is therefore $\omega\Delta t = 2\pi\Delta t/T$, and the freezing approximation holds whenever $\Delta t\ll T$.
In the GRN model, $\omega$ grows with interaction strength $|w_{ij}|$, degradation rate $\gamma_i$, and sigmoid slope $\phi'$, so fast oscillatory circuits impose the tightest requirements on temporal resolution.

\subsection{Trajectory Normalization and Temporal Rescaling}

The continuous dynamics of the network were initially modeled using the following stochastic differential equation (SDE):

\[dx_i(t) = -\gamma x_i(t)dt + V \sigma\Big(\sum_j J \Theta_{i,j} x_j(t)\Big)dt + D dW(t)\]


where $t \in [0, T]$ represents the physical time domain, $\gamma$ is the intrinsic decay rate, $V$ is the amplitude multiplier of the activation function $\sigma(\cdot)$, $J$ is the intensity multiplier of the coupling weights $\Theta_{i,j}$ which is a signed binary matrix, $D$ denotes the noise intensity, and $dW(t)$ represents the standard Wiener process increment. Under deterministic conditions, the state variables of this system naturally saturate at a maximum amplitude of $V/\gamma$.

To standardize the generated trajectories for comparative analysis, the system was mapped into a dimensionless framework where the state variables saturate at unity and the time domain spans the unit interval $s \in [0, 1]$. We applied a linear spatial scaling $y_i(s) = (\gamma/V) x_i(t)$ and a temporal scaling $s = t/T$, which implies $dt = T ds$. In accordance with the scaling properties of Brownian motion, the stochastic increment was transformed as $dW(t) = \sqrt{T} dW(s)$.

Substituting these transformations into the original governing equation yields the fully normalized SDE:


\[dy_i(s) = -\tilde{\gamma} y_i(s)ds + \tilde{V} \sigma\Big(\sum_j \tilde{J}\Theta_{i,j} y_j(s)\Big)ds + \tilde{D} dW(s)\]

This exact homeomorphic mapping is achieved by updating the system parameters as follows:
\begin{itemize}
    \item Effective decay rate: $\tilde{\gamma} = \gamma T$
    \item Effective amplitude: $\tilde{V} = \gamma T$
    \item Rescaled coupling weights: $\tilde{\Theta}_{i,j} = \Theta_{i,j} \frac{V}{\gamma}$
    \item Rescaled noise intensity: $\tilde{D} = D \frac{\gamma \sqrt{T}}{V}$
\end{itemize}

By utilizing these rescaled parameters, numerical integration of the normalized system (e.g., via the Euler-Maruyama method) guarantees structural equivalence to the original dynamics. This allows to map different systems with varying parameter regimes into a common normalized space, facilitating comparative analysis and inference while preserving the underlying dynamical behavior.


\subsection{Integrating prior knowledge}

Here we report in detail some possible prior functions that can be used to integrate prior knowledge about the network structure in the inference procedure.

\subsubsection{Presence of feedback loops}

To encourage a network structure with a high density of stabilizing motifs during inference, we define a structural prior that rewards the presence of negative feedback loops. This is achieved by defining a prior probability $P(G) \propto \exp(\lambda \cdot N_{neg})$, where $\lambda$ is a positive scaling factor and $N_{neg}$ is the count of 3-node negative feedback loops derived from the interaction matrix $W$ and its absolute counterpart $B = |A|$. The number of negative feedback loops corresponds to:
\[
N_{neg} = \frac{Tr(W^3)-Tr(B^3)}{6}
\]
where $Tr$ denotes the trace of a matrix. This prior effectively promotes network configurations that are enriched in negative feedback loops, which are known to contribute to the stability and robustness of gene regulatory networks.

\subsubsection{Sparsity prior}

Given a belived sparsity level $s^\ast$ of the network with an associated confidence level $\sigma_s\ast$, we can define a prior function that favors configurations with a number of non-zero weights close to the expected one. Such prior can be defined as:

\begin{equation}
    P(\mathbf{\Theta}|s^\ast,\sigma_s^\ast) = \exp\Big[-\frac{1}{2\sigma_s^\ast}\Big(\frac{\sum_{i,j}\mathbb{I}(\Theta_{ij}\neq 0)}{N^2}-s^\ast\Big)^2\Big]
\end{equation}

Where $\mathbb{I}(\Theta_{ij}\neq 0)$ is an indicator function that is equal to 1 if the weight $\Theta_{ij}$ is non-zero, and 0 otherwise. Here we chose the Gaussian form for the prior, but other forms can be used as well depending on the previous knowledge about the network structure.

\vspace{0.2 cm}

{\bf Acknowledgments}
F.Z. is supported by a Decreto Ministeriale n. 118 del 02/03/2023, M4C1 I. 4.1 from Piano Nazionale di Ripresa e Resilienza (PNRR). V.d.A. and M.D.D. acknowledge financial support from the EU funding within the MUR PNRR “National Center for HPC, BIG DATA AND QUANTUM COMPUTING” (Project no. CN00000013 CN1) and from the INFN grant “LINCOLN”. The authors acknowledge the generous support of CloudVeneto for the computing infrastructure.

\vspace{0.2 cm}
\vspace{0.2 cm}


{\bf Author contributions} 



\vspace{0.2 cm}



\renewcommand{\bibsection}{\section{References}}
\bibliographystyle{plainnat}
\bibliography{biblio}

\end{document}